\documentclass[crop=false]{standalone}
% --- ПРЕАМБУЛА ДЛЯ ПОДДЕРЖКИ РУССКОГО ЯЗЫКА И МАТЕМАТИКИ ---
\usepackage[T2A]{fontenc}
\usepackage[utf8]{inputenc}
\usepackage[russian]{babel}
\usepackage{amsmath}
\usepackage{geometry}
\usepackage{amssymb}
\usepackage{graphicx}
\usepackage{float}
\usepackage{import}
\geometry{a4paper, top=2cm, bottom=2cm, left=2cm, right=2cm}
\newcommand{\sidecomment}[1]{\marginpar{\raggedright\scriptsize\itshape #1}}
\newcommand{\iu}{\mathrm{i}}

\newcommand{\Def}{%
    \text{Def}%
    \hspace{-0.85em}% Сдвиг назад ровно на середину (подберите под шрифт)
    \raisebox{0.1ex}{/}% Черта
    \hspace{1em}% Возврат к концу слова + небольшой отступ
}

\renewcommand{\Re}{\operatorname{Re}}
\renewcommand{\Im}{\operatorname{Im}}

\ifstandalone
    \newcommand{\lecturebreak}{} % В отдельном файле лекции разрыв не нужен
\else
    \newcommand{\lecturebreak}{\clearpage} % В полной книге - начинаем с новой страницы
\fi

\newcounter{lecturenumber} % Создаем счетчик для нумерации лекций
\newcommand{\lecture}[2]{
    \lecturebreak
    \refstepcounter{lecturenumber} % Увеличиваем счетчик на 1
    \par\vspace{2em}\noindent % Отступ сверху и убираем абзацный отступ
    {\Large\bfseries % Делаем заголовок крупным и жирным
    Лекция \thelecturenumber: #1 % Выводим "Лекция N: Тема"
    \hfill % Растягиваем пробел, чтобы дата была справа
    \large\textit{#2}} % Выводим дату курсивом
    \par\vspace{1em} % Небольшой отступ снизу
    \addcontentsline{toc}{section}{Лекция \thelecturenumber: #1} % Добавляем в оглавление
    \par
}

\newcommand{\TwoCols}[2]{%
  \par\noindent % Начать с новой строки без отступа
  \begin{minipage}[t]{0.48\textwidth}
    #1
  \end{minipage}% <-- Важный '%' для удаления пробела
  \hfill % Растягивающийся пробел между колонками
  \begin{minipage}[t]{0.48\textwidth}
    #2
  \end{minipage}% <-- Важный '%' для удаления пробела
  \par%
}

\pagestyle{plain} % Включаем нумерацию страниц\

\begin{document}
\lecture{}{17.02}
\subsection{Некоторые свойства перемножения кватернионов}
\[a_1 \circ a_2 = a_1 a_2\]
\[a_1 \circ \bar a_2 = a_1 \bar a_2\]
\[a_1 \circ (a_2 + \bar a_2) = a_1 a_2 + a_1 \bar a_2 \]
\[\bar a_1 \circ \bar a_2 = - (\bar a_1, \bar a_2) + (\bar a_1 \times \bar a_2)\]
Сопряжённое:
\[\Lambda = \lambda + \bar \lambda \qquad \bar \Lambda = \lambda - \bar \lambda\]
\[\overline{q_1 \circ  q_2} = \bar q_2 \circ \bar q_1\]
\begin{proof}
\[\begin{aligned}
\overline{q_1 \circ  q_2} = \overline{(a_1 + \vec a_1) \circ (a_2 + \vec a_2)} &= \overline{a_1 a_2 - (\vec a_1, \vec a_2) + a_1 \vec a_2 + a_2 \vec a_1 + \vec a_1 \times \vec a_2} \\
& = a_1 a_1 - (a_1, a_2) - a_1 \vec a_2 - a_2 \vec a_1 - \vec a_1 \times \vec a_2
\end{aligned}
\]
\[\overline{(a_2 + \vec a_2)\circ (a_1 + \vec a_1)} = (a_2 - a\vec a_2) \circ (a_1 - \vec a_1) = a_1 a_2 - (\vec a_1,\vec a_2 ) - a_1 \vec a_2 - a_2 \vec a_1 + \vec a_2 \times \vec a_1\]
\end{proof}
Следствие: \[\overline{q \circ \bar q} = \bar{\bar q} \circ \bar q = q \circ \bar q = q_1^2 + q_2^2 + q_3^2 + q_4^2 = \|q\|\]
То есть векторная часть такого кватерниона нуль
\[|q| = \sqrt{\| q\| }= \sqrt{ q_1^2 + q_2^2 + q_3^2 + q_4^2 }\]
\[q^{-1} = \frac{\bar q}{\| q \|} \iff q^{-1} \circ q = \frac{\bar q \circ q}{\|q \|} = 1\]
\subsection{Описание поворотов}
\Def Пространство кватернионов обозначается $\mathbb H$
\[\psi : \mathbb H \to \mathbb H \qquad \psi_\Lambda(q) = \Lambda \circ q \circ \Lambda^{-1}\]
\[\bar \Lambda = - \Lambda \iff \Lambda = 0 + \vec \lambda\]
\[q = 0 + \vec a \qquad \psi(\vec a) = \Lambda \circ \vec a \circ \Lambda^{-1} \]
\[\overline{\psi(\vec a)} = \overline{\Lambda \circ \vec a \circ \Lambda^{-1}} = \bar \Lambda^{-1} \circ (- \vec a) \circ \bar \Lambda = \frac{1}{\| \Lambda \|} \bar {\bar \Lambda} \circ (- \vec a) \circ \bar \Lambda = - \Lambda \circ \vec a \circ \Lambda^{-1}\]
То есть вектор переходит в вектор
\[
\begin{aligned}
\vec a_1 \circ \vec a_2 \stackrel{?}{=}(\psi_\Lambda(\vec a_1) \circ \psi_\Lambda(\vec a_2)) &= \Lambda \circ \vec a_1 \circ \Lambda^{-1} \circ \Lambda \circ \vec a_2 \circ \Lambda^{-1}  = \Lambda \circ \vec a_1 \circ \vec a_2 \circ \Lambda^{-1} \\ 
&= \psi_\Lambda(\vec a_1 \circ \vec a_2) = \psi_\Lambda( -(\vec a_1, \vec a_2) + (\vec a_1 \times \vec a_2)) = \psi_\Lambda (- \vec a_1, \vec a_2) + \psi_\Lambda( \vec a_1 \times \vec a_2) 
\end{aligned}\]
\[\psi_\Lambda(\vec a_1) \circ \psi_\Lambda(\vec a_2) = - \underbracket{(\psi_\Lambda(\vec a_1), \psi_\Lambda(\vec a_2))} + (\psi_\Lambda(\vec a_1) \times \psi_\Lambda(\vec a_2)) = - \underbracket{\psi_\Lambda(a_1, a_2)} + \psi_\Lambda(\vec a_1 \times \vec a_2) \]
То есть $\psi$ сохраняет скалярное произведение

\subsection*{}
\keypoint 
\[
\begin{aligned}
\psi_\Lambda (\bar \lambda) = \Lambda \circ \lambda^{-1} \circ \Lambda^{-1}  &= (\lambda + \vec \lambda) \circ \vec \lambda \circ \frac{\lambda - \vec \lambda}{\| \Lambda \|} = (\lambda \circ \vec \lambda + \vec \lambda \circ \vec \lambda) \circ \frac{\lambda - \vec \lambda}{\| \Lambda \|} \\ 
&= \frac{1}{\| \Lambda \|} ( \lambda \circ \vec \lambda \circ \lambda - \lambda \circ \vec \lambda \circ \lambda + \vec \lambda \circ \vec \lambda \circ \lambda - \vec \lambda \circ \vec \lambda \circ \vec \lambda) = \frac{1}{\| \Lambda \|} (\lambda^2 \circ \vec \lambda - \vec \lambda \circ \vec \lambda \circ \vec \lambda) \\
&=  \frac{\lambda^2 + (\vec \lambda, \vec \lambda)}{\| \Lambda \|} \circ \vec \lambda = \vec \lambda
\end{aligned}
\]

\subsection{Кватернион поворота}
Вдоль оси $\vec u$ на угол $\varphi$
%есть фото доски
\[\Lambda = \cos \frac{\varphi}{2} + \sin \frac{\varphi}{2} \cdot \vec u\]
\[\vec e_z = \vec u \qquad \vec e_u = \frac{\bar u \times \bar r_0}{| \bar u \times \bar r_0|}\]
\[\bar r_0 = (\bar u ,\bar r_0) + \frac{\bar u \times \bar r_0 }{| \bar u \times \bar r_0 |} \times \bar u\]
\[\bar r = (\bar u, \bar r_0) + \cos \varphi (\bar u \times \bar r_0) \times \bar u + \sin \varphi (\bar u \times \bar r_0) \]
\[(\bar a\times \bar b) \times \vec \bar c = \bar b \times (\bar a \times \bar c) - \bar c \times ( \bar a \times \bar b)\]
Формула Родриго-Гамильтона:

\[\boxed{\bar r= \bar u (\bar u , \bar r_0) + \sin \varphi (\bar u \times \bar r_0) + \cos \varphi (\bar r_0 - \bar u (\bar u, \bar r_0))}\]
\[\begin{aligned}
    \bar r &= \Lambda \circ \bar r_0 \circ \Lambda^{-1} = (\cos \frac{\varphi}{2} + \sin \frac{\varphi}{2} \cdot \vec u) \circ \bar r_0 \circ (\cos \frac{\varphi}{2} - \sin \frac{\varphi}{2} \cdot \vec u) \\
    &= \left( \cos \frac{\varphi}{2} \bar r_0 + \sin \frac{\varphi}{2} \cdot \vec u \circ \bar r_0\right) \circ \left( \cos \frac{\varphi}{2} + \sin \frac{\varphi}{2} \cdot \vec u \right) \\
    &= \cos^2 \frac{\varphi}{2} \bar r_0 - \sin \frac{\varphi}{2} \cos \frac{\varphi}{2} ( - (\bar r_0 , \bar u) + (\bar r_0 \times \bar u)) - \sin \frac{\varphi}{2} \cos \frac{\varphi}{2} (\bar u, \bar r_0)+  \\
    &+ \sin \frac{\varphi}{2} \cos \frac{\varphi}{2} \bar u \times \bar r_0 - \sin^2 \frac{\varphi}{2} ( - (\bar u , \bar r_0) \bar u+  (\bar u \times \bar r_0) \circ \bar u) \\
    &= \cos^2 \frac{\varphi}{2} \bar r_0 + 2 \sin \frac{\varphi}{2} \cos \frac{\varphi}{2} (\bar u \times \bar r_0) + \sin^2 \frac{\varphi}{2}(\bar u , \bar r_0) \bar u - \sin^2 \frac{\varphi}{2} (\bar u \times \bar r_0) \times \bar u \\
    &= \cos^2 \frac{\varphi}{2} r_0 + \sin^2 \frac{\varphi}{2} (\bar u , \bar r_0) \bar u + \sin \varphi (\bar u \times \bar r_0) - \sin^2 \frac{\varphi}{2} (\bar u \times \bar r_0) \times \bar u
 \end{aligned}
    \]
\end{document}